% 已完成的工作

\subsection{工具链与编译脚本}

\begin{itemize}
    \item clang 汇编裸机代码（无 \texttt{-lSystem}）
    \item ld.lld 链接，指定加载地址
    \item 生成 ELF 镜像，QEMU 直接加载执行
    \item 固定 QEMU 参数：\texttt{-machine virt -cpu cortex-a72 -nographic}
\end{itemize}

编译运行脚本：

\begin{verbatim}
#!/bin/zsh
set -e
mkdir -p bin
/opt/homebrew/opt/llvm/bin/clang --target=aarch64-none-elf \
    -c start.s -o bin/start.o
/opt/homebrew/opt/lld/bin/ld.lld -Ttext=0x40000000 -e _start \
    -o bin/kernel.elf bin/start.o
qemu-system-aarch64 -machine virt -cpu cortex-a72 -nographic \
    -kernel bin/kernel.elf
\end{verbatim}

按 \texttt{Ctrl-A X} 退出 QEMU。

\subsection{启动层}

\begin{itemize}
    \item \texttt{\_start} 裸机入口，设置初始栈指针
    \item PL011 UART MMIO 输出，配合状态寄存器轮询
    \item 输出欢迎字符串并进入 halt 循环
\end{itemize}

核心启动汇编片段：

\begin{verbatim}
_start:
    ldr x0, =0x09000000
    ldr x1, =msg
print_loop:
    ldrb w2, [x1], #1
    cbz w2, halt
wait_uart:
    ldr w3, [x0, #0x18]
    tbz w3, #5, write_char
    b wait_uart
write_char:
    str w2, [x0]
    b print_loop
halt:
    wfe
    b halt
\end{verbatim}
